\chapter{Wnioski i kierunki dalszych badań}

\section{Podsumowanie wyników}

W niniejszej pracy przedstawiliśmy kompleksowe podejście do algorytmów uczenia ze wzmocnieniem w markowskich procesach decyzyjnych z dyskontowaniem quasi-hiperbolicznym. Główne osiągnięcia obejmują:

\begin{itemize}
\item Opracowanie algorytmu oceny polityki wolnego od modelu
\item Stworzenie algorytmu QH Q-learning z gwarancjami zbieżności
\item Zastosowanie algorytmów do modelu gromadzenia zapasów
\item Demonstrację niespójności czasowej optymalnych polityk
\end{itemize}

\section{Wkład naukowy}

Niniejsza praca wnosi istotny wkład do teorii i praktyki uczenia ze wzmocnieniem poprzez:
\begin{enumerate}
\item Rozszerzenie klasycznych algorytmów na przypadek dyskontowania quasi-hiperbolicznego
\item Udowodnienie zbieżności proponowanych metod
\item Pokazanie praktycznej stosowalności w problemach gospodarczych
\end{enumerate}

\section{Ograniczenia}

Przedstawione podejście ma następujące ograniczenia:
\begin{itemize}
\item Ograniczenie do przypadku decydenta zobowiązanego
\item Założenie o skończonych przestrzeniach stanów i akcji
\item Brak rozważenia przypadków częściowo obserwowalnych
\end{itemize}

\section{Kierunki dalszych badań}

Przyszłe badania mogą obejmować:
\begin{itemize}
\item Rozszerzenie na częściowo obserwowalne MDP (POMDP)
\item Systemy wieloagentowe z preferencjami czasowo-niespójnymi
\item Ciągłe przestrzenie stanów i akcji
\item Empiryczną walidację z udziałem ludzi
\item Estymację parametrów skłonności specyficznych dla jednostek
\end{itemize}

\section{Znaczenie praktyczne}

Wyniki pracy mają istotne znaczenie dla praktycznych zastosowań w:
\begin{itemize}
\item Finansach behawioralnych
\item Systemach rekomendacyjnych
\item Robotyce długoterminowej
\item Dynamicznym ustalaniu cen
\end{itemize}